% 来源：docs/04-maximum-likelihood.md
% 仅建立章节框架；正文、公式、图表与习题内容待后续迁移。
\chapter{最大似然与泛化}
\label{chap:mle}

\section*{本章要点}
% TODO：迁移原章“先抓住这几件事”的要点。
\contentplaceholder

\section{用 KL 散度比较两个分布}
\label{sec:mle-kl}
% TODO：补充本节内容。
\contentplaceholder

\subsection{离散与连续情形的 KL 定义}
\label{subsec:mle-definition}
% TODO：补充本节内容。
\contentplaceholder

\subsection{支持集条件与有限性}
\label{subsec:mle-support}
% TODO：补充本节内容。
\contentplaceholder

\subsection{KL 非负性的推导与取等条件}
\label{subsec:mle-nonnegativity}
% TODO：补充本节内容。
\contentplaceholder

\subsection{熵、交叉熵与 KL 的关系}
\label{subsec:mle-entropy}
% TODO：补充本节内容。
\contentplaceholder

\section{从分布拟合得到最大似然}
\label{sec:mle-distribution-fitting}
% TODO：补充本节内容。
\contentplaceholder

\subsection{总体正向 KL 与期望对数似然}
\label{subsec:mle-population}
% TODO：补充本节内容。
\contentplaceholder

\subsection{正向与反向 KL 的模式拟合倾向}
\label{subsec:mle-direction}
% TODO：补充本节内容。
\contentplaceholder

\subsection{独立同分布样本与经验对数似然}
\label{subsec:mle-empirical}
% TODO：补充本节内容。
\contentplaceholder

\subsection{似然、对数似然与平均负对数似然}
\label{subsec:mle-nll}
% TODO：补充本节内容。
\contentplaceholder

\section{两个可手算的例子}
\label{sec:mle-examples}
% TODO：补充本节内容。
\contentplaceholder

\subsection{硬币：频率就是伯努利参数的 MLE}
\label{subsec:mle-bernoulli}
% TODO：补充本节内容。
\contentplaceholder

\subsection{高斯均值：已知方差下得到样本平均}
\label{subsec:mle-gaussian}
% TODO：补充本节内容。
\contentplaceholder

\section{复杂模型：分解目标，再用梯度优化}
\label{sec:mle-optimization}
% TODO：补充本节内容。
\contentplaceholder

\subsection{联合对数似然的局部分解}
\label{subsec:mle-local}
% TODO：补充本节内容。
\contentplaceholder

\subsection{离散条件概率表的局部估计}
\label{subsec:mle-cpt}
% TODO：补充本节内容。
\contentplaceholder

\subsection{梯度上升、梯度下降与非凸优化}
\label{subsec:mle-gradient}
% TODO：补充本节内容。
\contentplaceholder

\subsection{小批量梯度估计及其条件}
\label{subsec:mle-minibatch}
% TODO：补充本节内容。
\contentplaceholder

\section{样本平均为何有用，又为何会过拟合？}
\label{sec:mle-generalization}
% TODO：补充本节内容。
\contentplaceholder

\subsection{蒙特卡洛平均的期望与方差}
\label{subsec:mle-monte-carlo}
% TODO：补充本节内容。
\contentplaceholder

\subsection{大数定律与数据依赖的参数估计}
\label{subsec:mle-data-dependent}
% TODO：补充本节内容。
\contentplaceholder

\subsection{总体风险、经验风险与验证风险}
\label{subsec:mle-risk}
% TODO：补充本节内容。
\contentplaceholder

\subsection{欠拟合、过拟合与偏差—方差视角}
\label{subsec:mle-fitting}
% TODO：补充本节内容。
\contentplaceholder

\subsection{模型容量与泛化界的适用条件}
\label{subsec:mle-bounds}
% TODO：补充本节内容。
\contentplaceholder

\subsection{训练似然、验证似然与生成质量}
\label{subsec:mle-evaluation}
% TODO：补充本节内容。
\contentplaceholder

\section{自检}
\label{sec:mle-review}
% TODO：补充本节内容。
\contentplaceholder

\section{来源与延伸阅读}
\label{sec:mle-reading}
% TODO：补充本节内容。
\contentplaceholder
